\documentclass[11pt]{article}
\usepackage{amsmath,amssymb,amsthm,algorithm,algorithmic}

\newtheorem{theorem}{Theorem}
\newtheorem{lemma}{Lemma}

\title{Problem 80: Square Root Digital Expansion}
\author{}
\date{}

\begin{document}
\maketitle

\section{Problem Statement}

For the first one hundred natural numbers, find the total of the digital sums of the first one hundred decimal digits for all irrational square roots.

\section{Mathematical Foundation}

\begin{theorem}[Irrationality Criterion]
For a positive integer $n$, $\sqrt{n}$ is irrational if and only if $n$ is not a perfect square.
\end{theorem}

\begin{proof}
($\Leftarrow$) Suppose $n$ is not a perfect square and $\sqrt{n} = a/b$ with $\gcd(a,b) = 1$. Then $nb^2 = a^2$. Let $p$ be a prime dividing $n$ to an odd power (such $p$ exists since $n$ is not a perfect square). Then $v_p(nb^2) = v_p(n) + 2v_p(b)$ is odd while $v_p(a^2) = 2v_p(a)$ is even, a contradiction.

($\Rightarrow$) If $n = k^2$, then $\sqrt{n} = k \in \mathbb{Q}$.
\end{proof}

\begin{lemma}[Perfect Squares]
The perfect squares in $\{1, \ldots, 100\}$ are $\{1, 4, 9, 16, 25, 36, 49, 64, 81, 100\}$ (10~values). We compute for $90$ non-square values.
\end{lemma}

\begin{proof}
$\lfloor\sqrt{100}\rfloor = 10$, giving exactly $10$ perfect squares $1^2, 2^2, \ldots, 10^2$.
\end{proof}

\begin{theorem}[High-Precision Extraction]
To obtain the first 100 decimal digits of $\sqrt{n}$ (non-square $n \leq 100$), compute $r = \lfloor\sqrt{n \cdot 10^{198}}\rfloor$ and take the first 100 digits of $r$.
\end{theorem}

\begin{proof}
We have $r = \lfloor 10^{99}\sqrt{n}\rfloor$ since $10^{198} = (10^{99})^2$. For $n \leq 100$, $\sqrt{n} < 11$, so $r < 11 \times 10^{99}$ has at most 101 digits. The first 100 digits give 100 significant decimal digits of $\sqrt{n}$.
\end{proof}

\begin{theorem}[Newton's Method for Integer Square Roots]
Given positive integer $N$, the sequence $x_0 = N$, $x_{k+1} = \lfloor(x_k + \lfloor N/x_k\rfloor)/2\rfloor$ converges to $\lfloor\sqrt{N}\rfloor$ with quadratic convergence. Termination: $x_{k+1} \geq x_k$ implies $x_k = \lfloor\sqrt{N}\rfloor$.
\end{theorem}

\begin{proof}
Let $s = \lfloor\sqrt{N}\rfloor$.

\textbf{(i)} For $x_k > s$: $N/x_k < x_k$, so $x_{k+1} \leq (x_k + N/x_k)/2 < x_k$ (strictly decreasing).

\textbf{(ii)} For $x_k > s$: by AM-GM, $(x_k + N/x_k)/2 \geq \sqrt{N} \geq s$, so $x_{k+1} \geq s$ (bounded below).

\textbf{(iii)} The sequence is strictly decreasing and bounded below by $s$, so it reaches $s$ in finite steps. At $x_k = s$: $\lfloor N/s\rfloor \in \{s, s+1\}$ (since $s^2 \leq N < (s+1)^2$), so $x_{k+1} = s \geq x_k$, and we terminate.

\textbf{(iv)} Quadratic convergence: $e_{k+1} \leq e_k^2/(2x_k)$ gives $O(\log\log N)$ iterations.
\end{proof}

\section{Algorithm}

\begin{algorithm}[H]
\caption{Sum of digital sums of first 100 digits of irrational $\sqrt{n}$ for $n = 1, \ldots, 100$}
\begin{algorithmic}[1]
\STATE $\text{total} \gets 0$
\FOR{$n = 1$ \TO $100$}
    \STATE $s \gets \text{isqrt}(n)$
    \IF{$s \cdot s = n$}
        \STATE \textbf{continue} \COMMENT{skip perfect squares}
    \ENDIF
    \STATE $N \gets n \cdot 10^{198}$
    \STATE $r \gets \text{isqrt}(N)$ \COMMENT{Newton's method}
    \STATE $\text{digits} \gets$ first 100 decimal digits of $r$
    \STATE $\text{total} \gets \text{total} + \sum \text{digits}$
\ENDFOR
\RETURN total
\end{algorithmic}
\end{algorithm}

\section{Complexity Analysis}

\textbf{Time:} For each of 90 non-squares, Newton's method on a $D \approx 200$ digit number takes $O(\log D)$ iterations, each costing $O(M(D))$ for division. With schoolbook arithmetic: $O(90 \cdot D^2 \cdot \log D)$.

\textbf{Space:} $O(D)$ per computation.

\section{Answer}

\[
\boxed{40886}
\]

\end{document}
